% !TeX root = ../main.tex
\chapter{Introducción}\label{cap:introduccion}

\section{Contexto y motivación}\label{sec:motivacion}

Presenta el problema, por qué es relevante y a quién afecta. Termina con
una o dos frases que expliquen qué aporta este trabajo.

Los acrónimos se escriben con \verb|\ac{...}|: la primera aparición de
\ac{ODS} se expande y las siguientes (\ac{ODS}) no.

\section{Objetivos}\label{sec:objetivos}

El objetivo general de este \ac{TFG} es \pendiente{redactar el objetivo general}.

Los objetivos específicos son:
\begin{enumerate}
  \item \textbf{O1.} Objetivo medible y verificable.
  \item \textbf{O2.} \ldots
  \item \textbf{O3.} \ldots
\end{enumerate}
El \cref{cap:conclusiones} revisa el grado de cumplimiento de cada uno.

\section{Metodología y planificación}\label{sec:planificacion}

Describe cómo se ha organizado el trabajo. El diagrama de la
\cref{fig:gantt} debe coincidir con los hitos (\textit{milestones}) del
tablero del proyecto en GitHub.

\begin{figure}[H]
  \centering
  \begin{ganttchart}[
      hgrid, vgrid={*1{draw=black!10}},
      x unit=0.3cm, y unit chart=0.6cm,
      title label font=\tiny,
      bar label font=\small,
      milestone label font=\small\itshape,
      bar/.append style={fill=blue!30},
    ]{1}{36}
    \gantttitle{Semanas}{36} \\
    \gantttitlelist{1,...,36}{1} \\
    \ganttbar{Arranque y propuesta}{1}{2} \\
    \ganttbar{Estado del arte}{2}{9} \\
    \ganttbar{Planteamiento}{6}{10} \\
    \ganttbar{Desarrollo}{9}{20} \\
    \ganttmilestone{Seguimiento}{15} \\
    \ganttbar{Resultados}{18}{25} \\
    \ganttbar{Redacción final}{22}{29} \\
    \ganttbar{Revisión y entrega}{29}{33} \\
    \ganttmilestone{Defensa}{36}
  \end{ganttchart}
  \caption{Planificación temporal del trabajo.}
  \label{fig:gantt}
\end{figure}

\section{Alineación con los Objetivos de Desarrollo Sostenible}\label{sec:ods}

Resume en un párrafo con qué \acs{ODS} se relaciona el trabajo. El
análisis detallado va en el \cref{anx:ods}.

\section{Estructura de la memoria}\label{sec:estructura}

El \cref{cap:estado-arte} revisa el estado del arte. El
\cref{cap:planteamiento} plantea el problema y la solución propuesta. El
\cref{cap:desarrollo} describe el desarrollo y el \cref{cap:resultados}
presenta los resultados. Por último, el \cref{cap:conclusiones} recoge las
conclusiones y las líneas de trabajo futuro.
